\documentclass[10pt,a4paper]{amsart}
\usepackage[margin=19mm]{geometry}
\usepackage{amsmath,amssymb,mathtools}
\usepackage[scr=rsfs,cal=euler]{mathalfa}
\usepackage{xfrac}
\usepackage[unicode,hidelinks,bookmarks=false]{hyperref}
\newcommand{\ie}{i.e.\ }
\newcommand{\cf}{cf.\ }
\newcommand{\resp}{respectively\ }
\newcommand{\Subsection}{\subsection}
\providecommand{\nobreakdash}{\nobreak\hbox{-}}
\setlength{\emergencystretch}{2em}
\multlinegap0pt
\usepackage{amssymb}
\usepackage{xcolor}
\newtheorem{lem}{Lemma}
\title{On segmentation by total variation type energies of Kobayashi-Warren-Carter type with fidelity}
\author{Yoshikazu Giga, Ayato Kubo, Hirotoshi Kuroda, Jun Okamoto, Koya Sakakibara}
\date{Source: arXiv:2408.04228v2 (CC BY 4.0)}
\begin{document}
\maketitle

\noindent\emph{Source and licence.} On segmentation by total variation type energies of Kobayashi-Warren-Carter type with fidelity. Version arXiv:2408.04228v2 (CC BY 4.0), distributed by arXiv under the Creative Commons Attribution 4.0 International licence, http://creativecommons.org/licenses/by/4.0/, as recorded in the arXiv licence metadata for this exact version and shown on the abstract page https://arxiv.org/abs/2408.04228v2. Full licence notice: \texttt{licenses/arxiv-2408.04228-CC-BY-4.0.txt}.\par\smallskip

\noindent\textit{Editorial scope.} Counted unit: the article's lem LPic (source0.tex lines 647--662). The article's complete original proof of it is delivered with it (source0.tex lines 664--781, one \texttt{proof} environment of 118 lines). No mathematics is written, repaired or re-derived: the statement and the proof are the article's own lines, in the article's own notation, and no retained line is substituted or re-flowed.  No further source-local context is needed: every cross-reference of every retained line resolves inside the two delivered spans.  One declared adaptation: exactly 3 ancillary strings inside the retained spans are omitted (image-inclusion commands and, where the source repeats a label, the repeated label definitions) and no other line differs from the source; the figure environments, with their placement, captions and labels, are kept, and the package records the omitted strings in fidelity receipts bound to the affected bodies.  No retained line contains a citation, so the wrapper carries no bibliography and no bibliography entry.  The retained lines contain the article's own diagram code, which the wrapper typesets with the standard tikz packages; no image file is delivered and the package declares no asset dependency.  Editorial formatting only, in addition: an AMS \texttt{amsart} preamble that restates the article's own declarations and macros used by the retained lines, namely the environments \texttt{lem} on separate counters, together with the standard packages those lines need.  The retained environments print in this document as Lemma 1 in order of appearance, whereas the article prints the same items with the numbering of its own full document.  The complete native source of the article as distributed in the arXiv e-print for this exact version is bundled unchanged as \texttt{sources/163507/source0.tex}.
\begin{lem} \label{LPic}
Assume that $K$ satisfies (K1) and that $g\in C[a,b]$.
 Let $U\in BV(a,b)$ be a minimizer of $TV_{Kg}$.
 Let $x_0\in(a,b)$ be a continuous point of $U$ and assume $U(x_0)>g(x_0)$ (resp.\ $U(x_0)<g(x_0)$).
 Then $U$ is constant in some interval $(\alpha,\beta)$ including $x_0$ and $U(x)>g(x)$ (resp.\ $U(x)<g(x)$) for $x\in(\alpha,\beta)$.
 Moreover, we can take $(\alpha,\beta)$ such that 
% 原稿 2024/4/15、2/6
one of following three cases occurs exclusively.
\begin{enumerate}
\item[(i)] $U(\alpha-0)<g(\alpha)<U(\alpha+0)$ (resp.\ $U(\alpha-0)>g(\alpha)>U(\alpha+0)$) and $U(\beta-0)=g(\beta)$,
\item[(i\hspace{-1pt}i)] $U(\alpha+0)=g(\alpha)$ and $U(\beta-0)>g(\beta)>U(\beta+0)$ (resp.\ $U(\beta-0)<g(\beta)<U(\beta+0)$),
\item[(i\hspace{-1pt}i\hspace{-1pt}i)] $U(\alpha+0)=g(\alpha)$ and $U(\beta-0)=g(\beta)$.
\end{enumerate}
In the case $\alpha=a$, we do not impose the condition $g(\alpha)>U(\alpha-0)$ (resp.\ $g(\alpha)<U(\alpha-0)$) and similarly, $g(\beta)>U(\beta+0)$ (resp.\ $g(\beta)<U(\beta+0)$) is not imposed for $\beta=b$ since $U(\alpha-0)$, $U(\beta+0)$ are undefined.
% The properties (i), (i\hspace{-1pt}i), (i\hspace{-1pt}i\hspace{-1pt}i) with opposite inequality hold if $U(x_0)<g(x_0)$.
\end{lem}

\begin{proof}
We shall only give a proof for the case $U(x_0)>g(x_0)$ since the argument for $U(x_0)<g(x_0)$ is symmetric.
 We consider the case that $U$ is continuous on $(\alpha,\beta)$ with $U>g$ on $(\alpha,\beta)$ and $U(\alpha+0)=g(\alpha)$, $U(\beta-0)=g(\beta)$.
 We shall claim that $U$ is a constant.
 If not, $\max_{[\alpha,\beta]}U>\min_{[\alpha,\beta]}U$.
 There would exist two points $x_1$ and $x_2$ in $[\alpha,\beta]$ such that $U(x_1)=\min_{[\alpha,\beta]}U$ and $U(x_2)=\max_{[\alpha,\beta]}U$ such that $U(x)\in(\min U,\max U)$ for $x$ between $x_1$ and $x_2$.
 We may assume $x_1<x_2$ with $U(x_1)<U(x_2)$ since the proof for the other case is symmetric.
 We first observe that $U$ must be a non-decreasing function on $[x_1,x_2]$.
 If not, there is either a local minimum or a local maximum of $U$ in $(x_1, x_2)$.
 In both cases, the argument is symmetric, we only discuss the case when there is a local maximum.
 In this case, there is $\gamma\in(x_1,x_2)$ and an open interval $I_\gamma=(y_1,y_2)\subset(x_1,x_2)$ containing $\gamma$ such that $\max_{\bar{I}_\gamma}U=U(\gamma)$ and $\max_{\partial I_\gamma}U<U(\gamma)$ with $U(y_1)=U(y_2)$.
 We set
\[
	v(x) := \left\{
\begin{array}{l}
	\max \left\{ U(y_1), U(x) - \varepsilon \right\}, \quad x \in I_\gamma \\
	U(x), \quad x \notin I_\gamma,
\end{array}
\right.
\]
where $\varepsilon$ is taken so that $U(x)-\varepsilon>\max\left(\max_{\partial I_\gamma}g, g(x)\right)$ for $x\in I_\gamma$.
 Since, $U>v>g$ on $I_\gamma$, we see that $\mathcal{F}(v)<\mathcal{F}(U)$.
 Evidently, $TV_K(v)=TV(v)\leq TV(U)=TV_K(U)$.
 This would contradict to the assumption that $U$ is a minimizer of $TV_{Kg}$.
 Hence, $U$ is monotone in $[x_1,x_2]$ satisfying $U>g$.
 Suppose that $U$ were not a constant function, there would exist another monotone continuous function $\tilde{U}$ ($\le U$) on $[x_1,x_2]$ such that $\tilde{U}(x_1)=U(x_1)$, $\tilde{U}(x_2)=U(x_2)$ keeping the property that $\tilde{U}>g$ on $(x_1,x_2)$.
 We set
\[
	v(x) = \left\{
\begin{array}{ll}
	\tilde{U}(x), & x \in (x_1, x_2), \\
	U(x), & x \notin (x_1, x_2).
\end{array}
\right.
\]
Since $U>v>g$ on $(x_1,x_2)$, we see that $\mathcal{F}(v)<\mathcal{F}(U)$.
 Clearly, $TV_K(v)=TV_K(U)$ so $U$ would not be a minimizer,
 This is a contradiction so we conclude that $U$ is a constant funtion on $[x_1,x_2]$ and (i\hspace{-1pt}i\hspace{-1pt}i) occurs.
 The proof shows that $U$ is constant in any interval $[x_1,x_2]$ provided that $U$ is continuous on $[x_1,x_2]$ and $U>g$ on $(x_1,x_2)$.

% 原稿 2024/4/15、3/6
Assume that there is a jump point $\alpha_0$ of $U$ with $\alpha_0<x_0$ with $U(\alpha_0+0)>g(\alpha_0)$.
 Then $U(\alpha_0+0)>U(\alpha_0-0)$.
 If not, $d=U(\alpha_0-0)-U(\alpha_0+0)>0$ and set
\[
	v(x) = \left\{
\begin{array}{lll}
	U(x) - d, & x \in (\alpha_0-\delta, \alpha_0) \\
	U(x), & x \notin (\alpha_0-\delta, \alpha_0).
\end{array}
\right.
\]
For a sufficiently small $\delta>0$, $v(x)>g(x)$; see Figure \ref{Fshi}.
\begin{figure}[tb]
\centering 

\caption{shift of a jump\label{Fshi}}
\end{figure}
By definition, $TV_K(v)\leq TV_K(U)$ and $\mathcal{F}(v)<\mathcal{F}(U)$.
 Thus $U$ is not a minimizer of $TV_{Kg}$.
 Similarly, if there is a jump point $\beta_0$ of $U$ with $x_0<\beta_0$, then $U(\beta_0-0)>U(\beta_0+0)$.

We shall prove that $U$ is continuous on $(x_0,\beta)$ provided that $U>g$ on $(\alpha_0,\beta)$.
 If not, there would exist a jump point $\beta_0$ of $U$ with $\alpha_0<x_0<\beta_0<\beta$ satisfying $U(x)>g(x)$ for $x\in(\alpha_0,\beta_0)$, $U(\alpha_0+0)>g(\alpha_0)$, $U(\beta_0-0)>g(\beta_0)$.
 We shall prove that such configuration does not occur.
 As we see in the previous paragraph, $U$ must be constant on $(\alpha_0,\beta_0)$ so that $U(\alpha_0+0)=U(\beta_0-0)$.
 We replace the value if $U$ a constant $M$ smaller than $U(\alpha_0+0)=U(\beta_0-0)$ but close to its value so that the new function $U_M$ still has a property that $U_M>g$.
 Since $U(\alpha_0-0)<U(\alpha_0+0)$, $U(\beta_0+0)<U(\beta_0-0)$, clearly $TV_K(U_M)\le TV_K(U)$ because of (K1).
 Moreover since $U\ge U_M$ and $U>U_M$ on $(\alpha_0,\beta_0)$ and $U_M>g$, $\mathcal{F}(U_M)<\mathcal{F}(U)$.
 Thus, $U$ is not a minimizer.
 Thus we observe that $\beta_0=\beta$ and we conclude that $U$ is a constant on $[x_0,\beta]$.
%\begin{figure}[tb]
%\centering 
%
%\caption{a way of truncation\label{Ftrun}}
%\end{figure}

Since there is a sequence of continuity point $x_j$ of $U$ converging to $\alpha_0$ as $j\to\infty$ keeping $x_j>\alpha_0$, we conclude that $U$ is a constant on $(x_j,\beta)$.
 This says that $U$ is constant on $(\alpha_0,\beta)$.
 To say that this corresponds to (i), it remains to prove that $U(\alpha-0)<g(\alpha)$.
 If not, $U(\alpha-0)\geq g(\alpha)$, then we take
\[
	v(x) = \left\{
\begin{array}{ll}
	\max \left( U(x) - d, \tilde{g}(x) \right), & x \in (\alpha_0, \alpha_0+\delta), \\
	U(x), & x \notin (\alpha_0, \alpha_0+\delta), \\
	U(x-0), & x = \alpha_0,
\end{array}
\right.
\]
where $d=U(\alpha_0+0)-U (\alpha_0-0)$.
 Here,
\[
	\tilde{g}(x) = \sup \left\{ g(y) \bigm|
	\alpha_0 \le y \le x \right\}
\]
which is continuous and non-decreasing.
 Moreover,
\[
	TV_K(v) \leq TV_K(U) + m\left(\tilde{g}(\delta)\right), \quad
	\mathcal{F}(v) \leq \mathcal{F}(U)
	- \delta \left(d-\tilde{g}(\delta)\right)^2,
\]
where $m(\sigma)=\sup\left\{ K(d+\tau)-K(d) \bigm| 0\leq\tau\leq\sigma\right\}$; see Figure \ref{Fmod}.
\begin{figure}[tb]
\centering 

\caption{modification of $U$\label{Fmod}}
\end{figure}
Thus
\[
	TV_{Kg}(v) \leq TV_{Kg}(U) + m\left(\tilde{g}(\delta)\right) - \delta\left(d-\tilde{g}(\delta)\right)^2.
\]
Since $\tilde{g}(\delta)\to0$ as $\delta\to0$, and $m(\sigma)\to0$ as $\sigma\to0$ by (K1), we conclude that for sufficiently small $\delta>0$, $TV_{Kg}(v)<TV_{Kg}(U)$.
 We now obtain (i).
 A symmetric argument yields (i\hspace{-1pt}i).
 The proof is now complete.
\end{proof}

\end{document}
